% Standalone Introduction and Related Work draft for:
% European Journal of Radiology Artificial Intelligence.
%
% The journal permits LaTeX submissions and limits Research Papers to 4,000
% words and 40 references. This file keeps the bibliography inline so that it
% can be compiled without a separate .bib file. When integrating it into the
% full manuscript, retain the sections and move the bibliography as needed.
%
% IMPORTANT: Elsevier currently requires a declaration of generative-AI use in
% manuscript preparation. The authors should add and personally verify that
% declaration in the final manuscript, before the reference list.

\documentclass[preprint,12pt]{elsarticle}

\usepackage[T1]{fontenc}
\usepackage{lmodern}
\usepackage{microtype}
\usepackage{amsmath}
\usepackage{amssymb}
\usepackage{booktabs}
\usepackage[hidelinks]{hyperref}
\usepackage{url}

\journal{European Journal of Radiology Artificial Intelligence}
\biboptions{sort&compress}

\begin{document}

\begin{frontmatter}

\title{VERA: Verifiable, Evidence-Grounded Regional Assembly for Faithful
Temporal Chest Radiograph Reporting}

\author{Author names omitted from this section draft}

\begin{abstract}
\textbf{Background:} Automated chest radiograph reporting systems can produce
clinically plausible text without ensuring that each statement is supported by
the image, particularly for anatomical location and interval change. This
disconnect limits clinical verification and may permit spatial or temporal
hallucinations.

\textbf{Objective:} To develop VERA (Verifiable, Evidence-grounded Regional
Assembly), a temporal reporting framework in which every report statement is a
traceable readout of localized, calibrated, and auditable image predictions.

\textbf{Methods:} VERA was developed using MIMIC-CXR studies and Chest ImaGenome
anatomical and longitudinal annotations with patient-level data separation. A
frozen biomedical vision--language encoder represents current and, when
available, prior radiographs; predicted bounding boxes define 29 anatomical
regions for masked attention pooling. A regional concept bottleneck maps 69
radiographic concepts to 14 observations through clinically specified,
non-negative connections. A temporal module predicts stable, improved, or
worsened status for valid region--observation pairs, and a deterministic
assembler applies calibrated assertion, hedging, omission, and temporal guards.
Evaluation covers anatomical localization, regional observation classification,
progression discrimination and calibration, concept and input-group
interventions, time reversal, claim provenance, and unsupported spatial or
temporal content.

% TODO: Replace the bracketed text below only after the final held-out and
% external evaluations have been frozen. Report the primary endpoint first and
% include confidence intervals where applicable.
\textbf{Results:} [Insert the principal held-out results for regional macro-F1,
progression macro-F1, calibration error, provenance coverage, spatial
unsupported-claim rate, and temporal-hallucination rate; include the most
informative baseline comparison and confidence intervals.]

\textbf{Conclusion:} VERA reframes chest radiograph report generation as
verifiable clinical evidence assembly rather than free-form diagnostic text
generation. Its constrained regional--temporal pathway is designed to make
evidence, uncertainty, and abstention explicit while preserving radiologist
oversight.
\end{abstract}

\end{frontmatter}

\section{Introduction}
\label{sec:introduction}

Chest radiography is among the most frequently performed imaging examinations,
and its interpretation is inherently longitudinal: radiologists routinely
compare a current examination with prior studies to distinguish genuine disease
progression from differences in inspiration, projection, positioning, and image
quality. Large public resources have accelerated automated chest radiograph
interpretation by pairing images with reports, global disease labels, anatomical
regions, and, more recently, localized comparison relations
\cite{johnson2019mimic,irvin2019chexpert,wu2021chestimagenome}. Nevertheless,
most automated systems still treat each examination as an isolated image or
produce a single image-level prediction. This formulation obscures the two
questions that are central to clinical verification: \emph{where} is the
abnormality, and \emph{what visual evidence changed relative to the prior
study}?

Automated radiology report generation offers a more expressive output, but
free-form generation introduces a safety-critical failure mode. A fluent report
may contain an unsupported finding, assign a correct finding to the wrong
anatomical location, omit a relevant abnormality, or describe an interval change
when no prior image is available. Memory-based, topic-guided, and retrieval-based
systems have improved language quality and clinical-label agreement, while
factual-reward and consistency objectives have reduced some contradictions
\cite{chen2020r2gen,nguyen2021accurate,endo2021repair,miura2021factual}.
However, these methods continue to express the diagnostic decision through a
generated or retrieved narrative. Consequently, a plausible sentence is not
necessarily a faithful account of the evidence that drove the image model, and
post-generation checking can only detect errors that are recoverable by the
chosen report labeler or information-extraction system.

Longitudinal chest radiograph models address part of this problem by explicitly
processing prior--current pairs. BioViL-T learns temporally informed
vision--language representations, CheXRelNet models anatomical relationships
between sequential studies, and TRACE couples temporal classification with
spatially grounded change descriptions
\cite{bannur2023biovilt,karwande2022chexrelnet,aranya2026trace}; more recent
direction-aware pretraining and temporally grounded benchmarks further show that
subtle states such as stability and resolution remain difficult under domain
shift \cite{hu2026protrans,prakash2026chextemporal}. These approaches establish
the value of temporal and spatial supervision, but they do not jointly enforce a
claim-level path from localized visual evidence, through a constrained disease
and progression decision, to every statement in the final report. In
particular, a generated temporal phrase can remain weakly coupled to the model's
actual progression computation.

We address this gap with \textbf{VERA} (\textbf{V}erifiable,
\textbf{E}vidence-grounded \textbf{R}egional \textbf{A}ssembly), an image-only
inference framework for faithful temporal chest radiograph reporting. Rather
than asking a language model to infer a diagnosis and subsequently explain it,
VERA represents the report as a deterministic readout of structured predictions.
A frozen BioViL-T encoder first provides a spatial feature grid for each image.
An anatomical detector localizes 29 predefined chest regions, and bbox-masked
attention pooling produces region-specific features whose pooling weights remain
available as within-region evidence. A regional classifier predicts 69 named
radiographic concepts and 14 CheXpert observations; its faithful concept path
uses a masked non-negative concept-to-disease mapping so that only clinically
mapped concepts can support each disease and concept interventions have a
defined monotonic effect. A temporal module then compares the current and prior
regional representations to predict \{stable, improved, worsened\} for each
valid region--disease pair. Finally, a report assembler applies calibrated
assert/hedge/omit decisions, attaches regional and temporal provenance, suppresses
all temporal language when no prior is present, and rejects text that cannot be
traced to the structured prediction table.

The framework deliberately separates supervision-time information from
deployment-time inputs. Reports and Chest ImaGenome scene graphs are used to
construct regional concept and progression targets, including through an
auditable deterministic report parser, but the deployed diagnostic path consumes
only the current image and, when available, a real prior image from the same
patient. Detector-predicted boxes are used for both training and inference to
avoid privileged anatomical localization at evaluation time; human-validated
boxes and temporal annotations are reserved for oracle analyses and external
audits. This separation is essential because training or evaluating on
report-derived labels is not equivalent to demonstrating that a progression
claim is supported by the images.

The principal contributions of this study are:

\begin{enumerate}
    \item a unified regional--temporal architecture that preserves the
    $29\times14$ anatomical disease structure from image encoding through
    progression assessment and report assembly, instead of collapsing evidence
    into a single global embedding or unconstrained narrative;

    \item a faithful regional concept bottleneck with a masked, non-negative
    concept-to-disease head, together with intervention and leakage tests that
    determine when concept-level ``why'' explanations are permitted;

    \item an explainable progression mechanism that combines localized temporal
    fusion, direction-aware current--prior reasoning, time-reversal consistency,
    and input-group interventions to expose whether a prediction is driven by
    the current image, the prior image, their explicit difference, or disease
    state changes; and

    \item a structured reporting and verification strategy in which every
    disease, location, confidence value, and temporal phrase is linked to an
    auditable model output, with hard abstention when the required evidence or
    comparison study is unavailable.
\end{enumerate}

This design shifts the objective from generating text that appears clinically
reasonable to producing claims whose computational and anatomical provenance can
be inspected. VERA is therefore intended as a decision-support and report-drafting
framework: it makes uncertainty and evidence explicit while preserving the
radiologist's responsibility for final interpretation.

\section{Related Work}
\label{sec:related_work}

\subsection{Chest radiograph datasets and vision--language representation
learning}

MIMIC-CXR and CheXpert made large-scale chest radiograph learning possible by
linking images to reports or automatically extracted observation labels, whereas
Chest ImaGenome added anatomy-centered scene graphs, 29 localized regions, and
longitudinal comparison relations \cite{johnson2019mimic,irvin2019chexpert,wu2021chestimagenome}.
These resources provide complementary supervision: global labels support robust
disease classification, reports provide rich but weak semantic signals, and
scene graphs permit regional and temporal reasoning. They also expose an
important methodological distinction between silver annotations extracted from
reports and gold annotations reviewed by humans; the former can support training
at scale, but final claims about localization and progression require independent
human-validated evaluation.

Medical vision--language pretraining uses naturally paired images and reports to
reduce dependence on manual labels. ConVIRT introduced bidirectional contrastive
learning for medical image--text pairs, GLoRIA aligned global and local visual
representations with report words, and BioViL strengthened the semantic modeling
of radiology language and phrase grounding
\cite{zhang2022convirt,huang2021gloria,boecking2022biovil}. BioViL-T subsequently
extended this paradigm to prior--current image sequences and temporal report
content \cite{bannur2023biovilt,karwande2022chexrelnet}. VERA uses a frozen
BioViL-T visual encoder for stable, reusable patch features, but differs in its
downstream objective: it preserves an explicit anatomical table and constrains
how concepts, diseases, and progression labels can influence the final report.
The pretrained representation is therefore an input to the reasoning pipeline,
not itself the explanation.

\subsection{Radiology report generation and factual consistency}

Early neural report generators adapted encoder--decoder architectures to the
repetitive and long-form structure of radiology text. R2Gen used relational
memory, topic-guided models introduced disease-oriented intermediate states, and
cross-modal or retrieval-based systems improved image--text alignment and
clinical-label accuracy \cite{chen2020r2gen,nguyen2021accurate,endo2021repair}.
Such models can produce fluent summaries and reduce repetitive drafting, yet
standard language metrics reward lexical similarity and may not penalize a
clinically consequential addition, omission, or contradiction.

Subsequent work introduced entity-level rewards, natural-language-inference
objectives, retrieval constraints, and specialized report labelers to improve or
measure factuality \cite{miura2021factual,ramesh2022priors,smit2020chexbert,khanna2023radgraph2}.
Notably, removing references to nonexistent priors reduces temporal hallucination
in training corpora, while CheXbert and RadGraph-style extraction make it
possible to compare generated text with structured clinical content. These are
valuable safeguards, but factual consistency with a target report or an
extracted label set does not guarantee faithfulness to the internal image
decision. VERA instead restricts report content before realization: the prose is
assembled from a calibrated table of allowed claims, and round-trip extraction
serves as a secondary verifier rather than the source of diagnostic truth.

\subsection{Longitudinal chest radiograph analysis}

Temporal chest radiograph analysis must separate disease evolution from changes
in acquisition geometry and anatomy. BioViL-T uses joint static--temporal
representation learning, CheXRelNet combines local and global anatomical
relations, and TRACE jointly predicts change descriptions and spatial grounding
\cite{bannur2023biovilt,karwande2022chexrelnet,aranya2026trace}. ProTrans further
formulates progression as a directional semantic transition with reversed-order
consistency, while CheXTemporal expands evaluation to localized five-state
transitions---new, worse, stable, improved, and resolved
\cite{hu2026protrans,prakash2026chextemporal}. Collectively, this literature
demonstrates that temporal direction and spatial correspondence must be modeled
together; however, most systems optimize a latent temporal representation or a
generated description without enforcing that each reported transition is a
readout of a constrained regional computation.

VERA retains the disease axis inherited from its regional classifier and predicts
progression for every supervised region--disease cell. Its temporal reasoning is
staged on a frozen regional model so that progression training cannot silently
alter the evidence used for single-image disease localization. In addition to
task metrics, the framework tests time reversal---stable should remain invariant,
whereas improved and worsened should exchange---and intervenes on explicit input
groups to determine whether the model truly uses the prior examination. This is
particularly important for detecting a degenerate temporal model that classifies
the current image accurately but ignores the comparison image.

\subsection{Explainability, concepts, and verifiable clinical output}

Post-hoc methods such as Grad-CAM and SHAP can localize influential pixels or
assign feature importance, but their explanations are computed after the model
has made its decision and do not, by themselves, impose a clinically meaningful
reasoning path \cite{selvaraju2017gradcam,lundberg2017shap}. Region annotations
and scene graphs add anatomical semantics, whereas concept bottleneck models make
human-readable variables part of the predictive computation and allow targeted
interventions \cite{wu2021chestimagenome,koh2020cbm,shin2023intervention}.
Concept-based approaches nevertheless remain vulnerable to incomplete concept
sets, noisy concept annotations, and bypass pathways in which the task head
ignores the purported explanation \cite{shin2023intervention,wang2026conceptcomplement}.

VERA combines anatomical grounding with a selective concept claim. Bbox-masked
pooling provides a direct ``where'' signal, while the principal concept-to-disease
path is masked and non-negative so that a disease cannot borrow an unrelated
concept and increasing a mapped concept cannot decrease its disease score.
Free-feature hybrid heads are treated as leakage-prone ablations rather than
faithful explanations, and concept text is exposed only for predictions that pass
concept-quality and intervention gates. For temporal predictions, the same
principle is extended from concept presence to concept change: supporting and
opposing current--prior concept contributions can be recorded in a progression
evidence ledger, whereas latent input-group effects are reported mechanistically
without being translated into unsupported clinical signs. This layered policy
allows the system to fall back from concept-level explanation to regional
grounding or abstention instead of presenting a plausible but unfaithful
rationale.

\begin{thebibliography}{00}

\bibitem{johnson2019mimic}
A.E.W. Johnson, T.J. Pollard, S.J. Berkowitz, N.R. Greenbaum,
M.P. Lungren, C.Y. Deng, R.G. Mark, S.J. Horng,
MIMIC-CXR, a de-identified publicly available database of chest radiographs
with free-text reports,
Sci. Data 6 (2019) 317.
\href{https://doi.org/10.1038/s41597-019-0322-0}{doi:10.1038/s41597-019-0322-0}.

\bibitem{irvin2019chexpert}
J. Irvin, P. Rajpurkar, M. Ko, et al.,
CheXpert: A large chest radiograph dataset with uncertainty labels and expert
comparison,
Proc. AAAI Conf. Artif. Intell. 33 (2019) 590--597.
\href{https://doi.org/10.1609/aaai.v33i01.3301590}{doi:10.1609/aaai.v33i01.3301590}.

\bibitem{wu2021chestimagenome}
J.T. Wu, N.N. Agu, I. Lourentzou, et al.,
Chest ImaGenome dataset for clinical reasoning,
Proc. NeurIPS Track Datasets Benchmarks 1 (2021).
\url{https://datasets-benchmarks-proceedings.neurips.cc/paper/2021/hash/17e62166fc8586dfa4d1bc0e1742c08b-Abstract-round2.html}.

\bibitem{zhang2022convirt}
Y. Zhang, H. Jiang, Y. Miura, C.D. Manning, C.P. Langlotz,
Contrastive learning of medical visual representations from paired images and
text,
Proc. Mach. Learn. Res. 182 (2022) 2--25.
\url{https://proceedings.mlr.press/v182/zhang22a.html}.

\bibitem{huang2021gloria}
S.-C. Huang, L. Shen, M.P. Lungren, S. Yeung,
GLoRIA: A multimodal global-local representation learning framework for
label-efficient medical image recognition,
Proc. IEEE/CVF Int. Conf. Comput. Vis. (2021) 3942--3951.
\url{https://openaccess.thecvf.com/content/ICCV2021/html/Huang_GLoRIA_A_Multimodal_Global-Local_Representation_Learning_Framework_for_Label-Efficient_Medical_ICCV_2021_paper.html}.

\bibitem{boecking2022biovil}
B. Boecking, N. Usuyama, S. Bannur, et al.,
Making the most of text semantics to improve biomedical vision--language
processing,
in: Computer Vision--ECCV 2022, Springer, 2022, pp. 1--21.
\href{https://doi.org/10.1007/978-3-031-20059-5_1}{doi:10.1007/978-3-031-20059-5\_1}.

\bibitem{bannur2023biovilt}
S. Bannur, S. Hyland, Q. Liu, et al.,
Learning to exploit temporal structure for biomedical vision--language
processing,
Proc. IEEE/CVF Conf. Comput. Vis. Pattern Recognit. (2023) 15016--15027.
\url{https://openaccess.thecvf.com/content/CVPR2023/html/Bannur_Learning_To_Exploit_Temporal_Structure_for_Biomedical_Vision-Language_Processing_CVPR_2023_paper.html}.

\bibitem{chen2020r2gen}
Z. Chen, Y. Song, T.-H. Chang, X. Wan,
Generating radiology reports via memory-driven Transformer,
Proc. 2020 Conf. Empir. Methods Nat. Lang. Process. (2020) 1439--1449.
\href{https://doi.org/10.18653/v1/2020.emnlp-main.112}{doi:10.18653/v1/2020.emnlp-main.112}.

\bibitem{nguyen2021accurate}
H. Nguyen, D. Nie, T. Badamdorj, et al.,
Automated generation of accurate and fluent medical X-ray reports,
Proc. 2021 Conf. Empir. Methods Nat. Lang. Process. (2021) 3552--3569.
\href{https://doi.org/10.18653/v1/2021.emnlp-main.288}{doi:10.18653/v1/2021.emnlp-main.288}.

\bibitem{endo2021repair}
M. Endo, R. Krishnan, V. Krishna, A.Y. Ng, P. Rajpurkar,
Retrieval-based chest X-ray report generation using a pre-trained contrastive
language-image model,
Proc. Mach. Learn. Res. 158 (2021) 209--219.
\url{https://proceedings.mlr.press/v158/endo21a.html}.

\bibitem{miura2021factual}
Y. Miura, Y. Zhang, E. Tsai, C.P. Langlotz, D. Jurafsky,
Improving factual completeness and consistency of image-to-text radiology
report generation,
Proc. NAACL-HLT (2021) 5288--5304.
\href{https://doi.org/10.18653/v1/2021.naacl-main.416}{doi:10.18653/v1/2021.naacl-main.416}.

\bibitem{ramesh2022priors}
V. Ramesh, N.A. Chi, P. Rajpurkar,
Improving radiology report generation systems by removing hallucinated
references to non-existent priors,
Proc. Mach. Learn. Res. 193 (2022) 456--473.
\url{https://proceedings.mlr.press/v193/ramesh22a.html}.

\bibitem{smit2020chexbert}
A. Smit, S. Jain, P. Rajpurkar, A. Pareek, A.Y. Ng, M.P. Lungren,
Combining automatic labelers and expert annotations for accurate radiology
report labeling using BERT,
Proc. 2020 Conf. Empir. Methods Nat. Lang. Process. (2020) 1500--1519.
\href{https://doi.org/10.18653/v1/2020.emnlp-main.117}{doi:10.18653/v1/2020.emnlp-main.117}.

\bibitem{khanna2023radgraph2}
S. Khanna, A. Dejl, K. Yoon, et al.,
RadGraph2: Modeling disease progression in radiology reports via hierarchical
information extraction,
Proc. Mach. Learn. Res. 219 (2023) 381--402.
\url{https://proceedings.mlr.press/v219/khanna23a.html}.

\bibitem{karwande2022chexrelnet}
G. Karwande, A. Mbakawe, J.T. Wu, L.A. Celi, M. Moradi, I. Lourentzou,
CheXRelNet: An anatomy-aware model for tracking longitudinal relationships
between chest X-rays,
arXiv preprint arXiv:2208.03873 (2022).
\href{https://doi.org/10.48550/arXiv.2208.03873}{doi:10.48550/arXiv.2208.03873}.

\bibitem{aranya2026trace}
O.F.M.R.R. Aranya, K. Desai,
TRACE: Temporal radiology with anatomical change explanation for grounded
X-ray report generation,
arXiv preprint arXiv:2602.02963 (2026).
\href{https://doi.org/10.48550/arXiv.2602.02963}{doi:10.48550/arXiv.2602.02963}.

\bibitem{hu2026protrans}
Z. Hu, Z. Yang, G. Wang, et al.,
Learning directional semantic transitions for longitudinal chest X-ray
analysis,
arXiv preprint arXiv:2606.15938 (2026).
\href{https://doi.org/10.48550/arXiv.2606.15938}{doi:10.48550/arXiv.2606.15938}.

\bibitem{prakash2026chextemporal}
E. Prakash, Y. Gao, C. Wang, et al.,
CheXTemporal: A dataset for temporally-grounded reasoning in chest radiography,
arXiv preprint arXiv:2605.11304 (2026).
\href{https://doi.org/10.48550/arXiv.2605.11304}{doi:10.48550/arXiv.2605.11304}.

\bibitem{selvaraju2017gradcam}
R.R. Selvaraju, M. Cogswell, A. Das, R. Vedantam, D. Parikh, D. Batra,
Grad-CAM: Visual explanations from deep networks via gradient-based
localization,
Proc. IEEE Int. Conf. Comput. Vis. (2017) 618--626.
\url{https://openaccess.thecvf.com/content_iccv_2017/html/Selvaraju_Grad-CAM_Visual_Explanations_ICCV_2017_paper.html}.

\bibitem{lundberg2017shap}
S.M. Lundberg, S.-I. Lee,
A unified approach to interpreting model predictions,
Adv. Neural Inf. Process. Syst. 30 (2017).
\url{https://proceedings.neurips.cc/paper/2017/hash/8a20a8621978632d76c43dfd28b67767-Abstract.html}.

\bibitem{koh2020cbm}
P.W. Koh, T. Nguyen, Y.S. Tang, et al.,
Concept bottleneck models,
Proc. Mach. Learn. Res. 119 (2020) 5338--5348.
\url{https://proceedings.mlr.press/v119/koh20a.html}.

\bibitem{shin2023intervention}
S. Shin, Y. Jo, S. Ahn, N. Lee,
A closer look at the intervention procedure of concept bottleneck models,
Proc. Mach. Learn. Res. 202 (2023) 31504--31520.
\url{https://proceedings.mlr.press/v202/shin23a.html}.

\bibitem{wang2026conceptcomplement}
H. Wang, J. Hou, S. He, S. Yang, H. Chen,
Concept complement bottleneck model for interpretable medical image diagnosis,
Proc. Mach. Learn. Res. 315 (2026) 342--359.
\url{https://proceedings.mlr.press/v315/wang26a.html}.

\end{thebibliography}

\end{document}
